% GENERATED FILE. Edit canonical lessons, then run npm run generate:books.
% canonical-source-hash: fnv1a64:2076fa8429b86a93
% canonical-lessons: SA-C471-mastiskah, SA-C471-kapalam, SA-C471-smasru, SA-C471-paksma, SA-C471-mundayati, SA-R471-review-464-471

\chapter{Head and Face}
\label{ch:sa-head-and-face}
\hlchaptermodality{471}

\begin{chapteropening}
\textbf{By the end of this chapter:} \emph{I can talk about the head, the beard and the eyelashes.}

Complete the last lesson of chapter 471: I can talk about the head, the beard and the eyelashes.
\end{chapteropening}

\section[\texorpdfstring{\sk{मस्तिष्कः} (mastiṣkaḥ)}{मस्तिष्कः (mastiṣkaḥ)}]{\sk{मस्तिष्कः} (mastiṣkaḥ) — the brain}

\label{lesson:SA-C471-mastiskah}



\begin{quote}
Before the new one: say the Sanskrit for the stomach, then the Sanskrit for breathes.
\end{quote}

\subsection*{You'll want to know: \sk{मस्तिष्कः}}
\textbf{\sk{मस्तिष्कः}} — \emph{mastiṣkaḥ} — \textquotedblleft{}the brain\textquotedblright{}.

\begin{quote}
\sk{मस्तिष्केण} \sk{नरः} \sk{चिन्तयति।}
\end{quote}

\begin{quote}
\emph{mastiṣkeṇa naraḥ cintayati.}
\end{quote}

\begin{quote}
A person thinks with the brain.
\end{quote}

A masculine noun.

\begin{cousinweb}
\textbf{\sk{मस्तिष्कः}} (\emph{mastiṣkaḥ}) MW gives \textquotedblleft{}the brain\textquotedblright{} from the Rigveda on; MW gives no derivation, so its origin is uncertain (MW).
\end{cousinweb}

\begin{grammarlens}[title={What you've built}]
One more word for describing what you lived and felt.
\end{grammarlens}

\subsection*{Your turn}
\begin{itemize}
\raggedright
  \item \emph{Say it:} \emph{mastiṣkaḥ}
  \item \emph{Say it:} \emph{mastiṣkaḥ}, once more
  \item \emph{Say it:} \emph{mastiṣkeṇa naraḥ cintayati.} — A person thinks with the brain.
  \item \emph{Recall:} say the Sanskrit for the stomach, then the Sanskrit for breathes, then say \emph{mastiṣkaḥ} again
\end{itemize}

\begin{tcolorbox}[breakable,colback=teal!4,colframe=teal!35!black,title={Before you move on}]
What does \sk{मस्तिष्कः} mean? (\textquotedblleft{}The brain\textquotedblright{}.) Say it once more.
\end{tcolorbox}
\section[\texorpdfstring{\sk{कपालम्} (kapālam)}{कपालम् (kapālam)}]{\sk{कपालम्} (kapālam) — the skull}

\label{lesson:SA-C471-kapalam}



\begin{quote}
Before the new one: say the Sanskrit for breathes, then the Sanskrit for the brain.
\end{quote}

\subsection*{You'll want to know: \sk{कपालम्}}
\textbf{\sk{कपालम्}} — \emph{kapālam} — \textquotedblleft{}the skull\textquotedblright{}.

\begin{quote}
\sk{कठिनं} \sk{कपालं} \sk{मस्तिष्कं} \sk{रक्षति।}
\end{quote}

\begin{quote}
\emph{kaṭhinaṃ kapālaṃ mastiṣkaṃ rakṣati.}
\end{quote}

\begin{quote}
The hard skull protects the brain.
\end{quote}

A neuter noun.

\begin{cousinweb}
\textbf{\sk{कपालम्}} (\emph{kapālam}) is traced to the root √kamp in the Uṇādi list that MW follows; its first senses are \textquotedblleft{}a cup, a dish\textquotedblright{}, and \textquotedblleft{}the skull, the cranium\textquotedblright{} comes after them (MW).
\end{cousinweb}

\begin{grammarlens}[title={What you've built}]
One more word for describing what you lived and felt.
\end{grammarlens}

\subsection*{Your turn}
\begin{itemize}
\raggedright
  \item \emph{Say it:} \emph{kapālam}
  \item \emph{Say it:} \emph{kapālam}, once more
  \item \emph{Say it:} \emph{kaṭhinaṃ kapālaṃ mastiṣkaṃ rakṣati.} — The hard skull protects the brain.
  \item \emph{Recall:} say the Sanskrit for breathes, then the Sanskrit for the brain, then say \emph{kapālam} again
\end{itemize}

\begin{tcolorbox}[breakable,colback=teal!4,colframe=teal!35!black,title={Before you move on}]
What does \sk{कपालम्} mean? (\textquotedblleft{}The skull\textquotedblright{}.) Say it once more.
\end{tcolorbox}
\section[\texorpdfstring{\sk{श्मश्रु} (śmaśru)}{श्मश्रु (śmaśru)}]{\sk{श्मश्रु} (śmaśru) — a beard}

\label{lesson:SA-C471-smasru}



\begin{quote}
Before the new one: say the Sanskrit for the brain, then the Sanskrit for the skull.
\end{quote}

\subsection*{You'll want to know: \sk{श्मश्रु}}
\textbf{\sk{श्मश्रु}} — \emph{śmaśru} — \textquotedblleft{}a beard\textquotedblright{}.

\begin{quote}
\sk{पितामहस्य} \sk{श्मश्रु} \sk{श्वेतम्} \sk{अस्ति।}
\end{quote}

\begin{quote}
\emph{pitāmahasya śmaśru śvetam asti.}
\end{quote}

\begin{quote}
Grandfather's beard is white.
\end{quote}

A neuter noun.

\begin{cousinweb}
\textbf{\sk{श्मश्रु}} (\emph{śmaśru}) MW calls it of unknown derivation, so its origin is uncertain; it gives \textquotedblleft{}the beard, esp. the moustache\textquotedblright{} (MW).
\end{cousinweb}

\begin{grammarlens}[title={What you've built}]
One more word for describing what you lived and felt.
\end{grammarlens}

\subsection*{Your turn}
\begin{itemize}
\raggedright
  \item \emph{Say it:} \emph{śmaśru}
  \item \emph{Say it:} \emph{śmaśru}, once more
  \item \emph{Say it:} \emph{pitāmahasya śmaśru śvetam asti.} — Grandfather's beard is white.
  \item \emph{Recall:} say the Sanskrit for the brain, then the Sanskrit for the skull, then say \emph{śmaśru} again
\end{itemize}

\begin{tcolorbox}[breakable,colback=teal!4,colframe=teal!35!black,title={Before you move on}]
What does \sk{श्मश्रु} mean? (\textquotedblleft{}A beard\textquotedblright{}.) Say it once more.
\end{tcolorbox}
\section[\texorpdfstring{\sk{पक्ष्म} (pakṣma)}{पक्ष्म (pakṣma)}]{\sk{पक्ष्म} (pakṣma) — an eyelash}

\label{lesson:SA-C471-paksma}



\begin{quote}
Before the new one: say the Sanskrit for the skull, then the Sanskrit for a beard.
\end{quote}

\subsection*{You'll want to know: \sk{पक्ष्म}}
\textbf{\sk{पक्ष्म}} — \emph{pakṣma} — \textquotedblleft{}an eyelash\textquotedblright{}.

\begin{quote}
\sk{मम} \sk{नेत्रे} \sk{पक्ष्म} \sk{अस्ति।}
\end{quote}

\begin{quote}
\emph{mama netre pakṣma asti.}
\end{quote}

\begin{quote}
There is an eyelash in my eye.
\end{quote}

A neuter noun.

\begin{cousinweb}
\textbf{\sk{पक्ष्म}} (\emph{pakṣma}) MW records \emph{pakṣman}, \textquotedblleft{}the eyelashes\textquotedblright{}; MW gives no derivation, so its origin is uncertain (MW).
\end{cousinweb}

\begin{grammarlens}[title={What you've built}]
One more word for describing what you lived and felt.
\end{grammarlens}

\subsection*{Your turn}
\begin{itemize}
\raggedright
  \item \emph{Say it:} \emph{pakṣma}
  \item \emph{Say it:} \emph{pakṣma}, once more
  \item \emph{Say it:} \emph{mama netre pakṣma asti.} — There is an eyelash in my eye.
  \item \emph{Recall:} say the Sanskrit for the skull, then the Sanskrit for a beard, then say \emph{pakṣma} again
\end{itemize}

\begin{tcolorbox}[breakable,colback=teal!4,colframe=teal!35!black,title={Before you move on}]
What does \sk{पक्ष्म} mean? (\textquotedblleft{}An eyelash\textquotedblright{}.) Say it once more.
\end{tcolorbox}
\section[\texorpdfstring{\sk{मुण्डयति} (muṇḍayati)}{मुण्डयति (muṇḍayati)}]{\sk{मुण्डयति} (muṇḍayati) — shaves}

\label{lesson:SA-C471-mundayati}



\begin{quote}
Before the new one: say the Sanskrit for a beard, then the Sanskrit for an eyelash.
\end{quote}

\subsection*{You'll want to know: \sk{मुण्डयति}}
\textbf{\sk{मुण्डयति}} — \emph{muṇḍayati} — \textquotedblleft{}shaves\textquotedblright{}.

\begin{quote}
\sk{नापितः} \sk{मम} \sk{पितुः} \sk{श्मश्रु} \sk{मुण्डयति।}
\end{quote}

\begin{quote}
\emph{nāpitaḥ mama pituḥ śmaśru muṇḍayati.}
\end{quote}

\begin{quote}
The barber shaves my father's beard.
\end{quote}

Past (he …): \sk{सः} \sk{शिरः} \sk{मुण्डितवान्}, \emph{saḥ śiraḥ muṇḍitavān}.

\begin{cousinweb}
\textbf{\sk{मुण्डयति}} (\emph{muṇḍayati}) is traced by MW to a root √muṇḍ, \textquotedblleft{}to shave\textquotedblright{}, which MW thinks probably made up to explain words such as \emph{muṇḍa}, \textquotedblleft{}shaven\textquotedblright{}; its history is not settled (MW).
\end{cousinweb}

\begin{grammarlens}[title={What you've built}]
One more word for describing what you lived and felt.
\end{grammarlens}

\subsection*{Your turn}
\begin{itemize}
\raggedright
  \item \emph{Say it:} \emph{muṇḍayati}
  \item \emph{Say it:} \emph{muṇḍayati}, once more
  \item \emph{Say it:} \emph{nāpitaḥ mama pituḥ śmaśru muṇḍayati.} — The barber shaves my father's beard.
  \item \emph{Recall:} say the Sanskrit for a beard, then the Sanskrit for an eyelash, then say \emph{muṇḍayati} again
\end{itemize}

\begin{tcolorbox}[breakable,colback=teal!4,colframe=teal!35!black,title={Before you move on}]
What does \sk{मुण्डयति} mean? (\textquotedblleft{}Shaves\textquotedblright{}.) Say it once more.
\end{tcolorbox}
\section[(dialogue)]{Looking back: from Talkers and Quiet Ones to Head and Face}

\label{lesson:SA-R471-review-464-471}



\begin{quote}
Say the Sanskrit for talkative and of few words.
\end{quote}

\subsection*{Your turn}
Say these aloud:

\begin{itemize}
\raggedright
  \item talkative, of few words, calls out to, reviles, prattle
  \item fickle, indifferent, is careless, attention, industrious
  \item mental vigour, is bewildered, madness, childish, understanding
  \item well-bred, mockery, impudent, flattery, affable
  \item greedy, a favour in return, a wrong done, reminds, grateful
  \item a sense, eyesight, touch, a smell, smells
  \item the lungs, the liver, a kidney, the stomach, breathes
  \item the brain, the skull, a beard, an eyelash, shaves
\end{itemize}

\begin{tcolorbox}[breakable,colback=teal!4,colframe=teal!35!black,title={Before you move on}]
Talkative? (\textbf{\sk{वाचालः}}, \emph{vācālaḥ}.) Of few words? (\textbf{\sk{मितभाषी}}, \emph{mitabhāṣī}.) Calls out to? (\textbf{\sk{सम्बोधयति}}, \emph{sambodhayati}.) Reviles? (\textbf{\sk{आक्रोशति}}, \emph{ākrośati}.) Prattle? (\textbf{\sk{प्रलापः}}, \emph{pralāpaḥ}.) Fickle? (\textbf{\sk{चञ्चलः}}, \emph{cañcalaḥ}.) Indifferent? (\textbf{\sk{उदासीनः}}, \emph{udāsīnaḥ}.) Is careless? (\textbf{\sk{प्रमाद्यति}}, \emph{pramādyati}.) Attention? (\textbf{\sk{अवधानम्}}, \emph{avadhānam}.) Industrious? (\textbf{\sk{उद्यमी}}, \emph{udyamī}.) Mental vigour? (\textbf{\sk{मेधा}}, \emph{medhā}.) Is bewildered? (\textbf{\sk{मुह्यति}}, \emph{muhyati}.) Madness? (\textbf{\sk{उन्मादः}}, \emph{unmādaḥ}.) Childish? (\textbf{\sk{बालिशः}}, \emph{bāliśaḥ}.) Understanding? (\textbf{\sk{बोधः}}, \emph{bodhaḥ}.) Well-bred? (\textbf{\sk{शिष्टः}}, \emph{śiṣṭaḥ}.) Mockery? (\textbf{\sk{उपहासः}}, \emph{upahāsaḥ}.) Impudent? (\textbf{\sk{धृष्टः}}, \emph{dhṛṣṭaḥ}.) Flattery? (\textbf{\sk{चाटुः}}, \emph{cāṭuḥ}.) Affable? (\textbf{\sk{प्रियंवदः}}, \emph{priyaṃvadaḥ}.) Greedy? (\textbf{\sk{लोलुपः}}, \emph{lolupaḥ}.) A favour in return? (\textbf{\sk{प्रत्युपकारः}}, \emph{pratyupakāraḥ}.) A wrong done? (\textbf{\sk{अपकारः}}, \emph{apakāraḥ}.) Reminds? (\textbf{\sk{स्मारयति}}, \emph{smārayati}.) Grateful? (\textbf{\sk{कृतवेदी}}, \emph{kṛtavedī}.) A sense? (\textbf{\sk{इन्द्रियम्}}, \emph{indriyam}.) Eyesight? (\textbf{\sk{दृष्टिः}}, \emph{dṛṣṭiḥ}.) Touch? (\textbf{\sk{स्पर्शः}}, \emph{sparśaḥ}.) A smell? (\textbf{\sk{गन्धः}}, \emph{gandhaḥ}.) Smells? (\textbf{\sk{जिघ्रति}}, \emph{jighrati}.) The lungs? (\textbf{\sk{फुप्फुसः}}, \emph{phupphusaḥ}.) The liver? (\textbf{\sk{यकृत्}}, \emph{yakṛt}.) A kidney? (\textbf{\sk{वृक्कः}}, \emph{vṛkkaḥ}.) The stomach? (\textbf{\sk{आमाशयः}}, \emph{āmāśayaḥ}.) Breathes? (\textbf{\sk{श्वसिति}}, \emph{śvasiti}.) The brain? (\textbf{\sk{मस्तिष्कः}}, \emph{mastiṣkaḥ}.) The skull? (\textbf{\sk{कपालम्}}, \emph{kapālam}.) A beard? (\textbf{\sk{श्मश्रु}}, \emph{śmaśru}.) An eyelash? (\textbf{\sk{पक्ष्म}}, \emph{pakṣma}.) Shaves? (\textbf{\sk{मुण्डयति}}, \emph{muṇḍayati}.)
\end{tcolorbox}
